% !TeX root = main.tex
\chapter{Модель мультимодального представлення прихованої загрози в мультимедійних об’єктах}
\label{chap:models_formalization}

У першому розділі встановлено, що ознаки прихованого навантаження можуть спостерігатися у розкодованому вмісті, частотному або багатомасштабному поданні, структурі контейнера та службових полях, тоді як прихований вплив установлюють за фактичною причинно пов’язаною реакцією під час контрольованого оброблення мультимедійного об’єкта. Проте окреме відхилення не доводить наявності аномального коду. Воно може бути наслідком допустимого перетворення, особливостей формату або неповноти спостереження.

Наукова суперечність полягає в багаторівневому прояві прихованої загрози за переважно ізольованого формування локальних оцінок. Для її розв’язання модель має спільно враховувати різнорідні ознаки, зберігати сильне свідчення окремого джерела, позначати відсутні спостереження, враховувати залежність ознак і не посилювати висновок понад наявні підстави. Просте об’єднання числових показників цього не забезпечує, оскільки приховує походження ознаки, умови її одержання та причину відсутності окремого спостереження. Для розв’язання зазначеної суперечності розроблено модель мультимодального представлення прихованої загрози. Її призначенням є узгоджене представлення різнорідних відомостей зі збереженням їхнього походження до моменту прийняття рішення.

Предметне тлумачення моделі ґрунтується на розмежуванні шести основних понять:

\begin{enumerate}[label=\arabic*),leftmargin=1.1cm,itemsep=0pt,parsep=0pt]
\item \textit{Аномалія}. Спостережуване відхилення від установленої структурної, статистичної або поведінкової норми. Окрема аномалія не встановлює наявності загрози.
\item \textit{Аномальний код}. Прихована послідовність даних, команда, фрагмент коду або структурне значення, що одночасно не відповідає заявленому призначенню мультимедійного об’єкта і створює умови для несанкціонованої дії, зокрема відновлення керувальних даних або небезпечної взаємодії з вебсистемою.
\item \textit{Прихована загроза}. Сценарій можливого небезпечного впливу, для якого встановлюється зв’язок між властивостями носія, операцією оброблення та очікуваним наслідком.
\item \textit{Шкідливий код}. Код, програмну семантику і шкідливу функцію якого підтверджено.
\item \textit{Приховане навантаження}. Дані, подані в елементах мультимедійного об’єкта за правилом прихованого розміщення та доступні для відновлення за відповідним правилом оброблення.
\item \begin{minipage}[t]{\linewidth}
\textit{Прихований вплив}. Фактичний спостережуваний наслідок опрацювання прихованого навантаження конкретним компонентом вебсистеми, причинний зв’язок якого з досліджуваним об’єктом установлено в межах доступних спостережень.
\end{minipage}
\end{enumerate}

Таке розмежування запобігає ототожненню виявленого відхилення з доведеним виконанням шкідливої функції.

Побудова моделі ґрунтується на п’яти припущеннях. Початковий мультимедійний об’єкт зберігається незмінним. Для кожного похідного представлення відоме правило його одержання. Відсутність представлення не тлумачиться як нульове значення інтегральної оцінки доказових підстав. Подію відносять до досліджуваного об’єкта лише за підтвердженого причинного зв’язку. Одержаний висновок характеризує стан об’єкта за наявних даних і чинних параметрів, але не гарантує його абсолютної безпечності.

Межу довіри проведено між збереженою початковою байтовою послідовністю та результатами її тлумачення. Послідовність $X_{mm}$ достовірно відтворює вхідні дані, однак сама по собі не підтверджує заявленого формату чи безпечності вмісту. Засоби структурного аналізу й розкодування, а також контрольоване середовище мають власні обмеження і можуть сформувати неповний результат. Тому разом з ознаками зберігають версії застосованих засобів, показники якості, правило одержання похідних даних і причину їхньої неповноти.

Числові результати можна порівнювати лише за сумісних схем ознак та однакових умов спостереження. Зміна засобу розкодування, правила приймання, послідовності оброблення або доступності журналу може змінити тлумачення ознаки без зміни початкового об’єкта. Тому повторна перевірка доповнює попередні спостереження, зберігаючи умови кожного запуску.

Наукова відмінність моделі полягає у типізованому поєднанні незмінного байтового подання, похідних представлень, локалізацій, маски доступності та відомостей про походження в межах однієї перевірки, а також у відокремленні аналітичного стану від технологічної дії.

Логіку моделі утворюють взаємозалежні формальні переходи від умов проходження мультимедійного об’єкта до його багаторівневого представлення та типізованих ознак. Узгодження цих ознак формує аналітичний стан з інтегральною оцінкою доказових підстав, упевненістю й поясненням, тоді як межі передавання результатів відокремлюють модельні висновки від методів та дії інформаційної технології.

Математичні залежності подано для одного мультимедійного об’єкта без наскрізного індексу. Під час опису вибірки або повторних запусків до відповідного позначення додають індекс $i$, наприклад $X_{mm,i}$ або $Z_i$, не змінюючи змісту величини.

\section{Умови здійснення атак на вебсистеми з використанням мультимедійних об’єктів}
\label{sec:attack_model}

У досліджуваному сценарії мультимедійний об’єкт виконує функцію носія прихованого навантаження та засобу його передавання через штатний або формально дозволений канал. Небезпека зумовлена відмінністю правил тлумачення на послідовних етапах оброблення. Об’єкт може пройти первинну перевірку, але його розкодовані значення, структуру контейнера, службові поля або часову організацію відео різні компоненти можуть тлумачити по-різному.

Повний перелік способів внесення прихованих даних наперед визначити неможливо. Тому модель описує не конкретний спосіб атаки, а зв’язок між початковим об’єктом, прихованим навантаженням і результатом цілеспрямованої підготовки носія. Цей зв’язок визначено формулою~\eqref{eq:attack_embedding}:
\begin{equation}
X_{mm}^{\prime}=E_A(X_{mm},P_A,\theta_A),
\label{eq:attack_embedding}
\end{equation}

де $X_{mm}$ позначає початковий мультимедійний об’єкт; $P_A$ є прихованим навантаженням; $\theta_A\in\mathrm T_A$ є конкретним кортежем параметрів підготовки з множини допустимих кортежів $\mathrm T_A$; $E_A$ позначає функцію внесення даних або цілеспрямованої зміни; $X_{mm}^{\prime}$ є підготовленим мультимедійним об’єктом.

Формула~\eqref{eq:attack_embedding} установлює межу подальшого розгляду. Вона описує результат цілеспрямованої підготовки, не визначаючи способу внесення даних і не передбачаючи обов’язкової зміни зовнішнього вигляду об’єкта. Подальший аналіз ґрунтується на властивостях $X_{mm}^{\prime}$, які можна встановити під час контрольованої перевірки. Надалі, якщо порівняння з носієм до цілеспрямованої підготовки не виконується, прийнятий до перевірки об’єкт $X_{mm}^{\prime}$ для стислості позначено $X_{mm}$.

Функція $E_A$, навантаження $P_A$, початковий і підготовлений об’єкти $X_{mm}$ та $X_{mm}^{\prime}$, а також величини $D$ і $S$ належать до контрольованого або еталонного сценарію підготовки носія. Під час аналізу невідомого вхідного об’єкта ці складники не є безпосередньо доступними ознаками. У такому разі операційний аналіз починається зі збереженої незмінної вхідної байтової послідовності та похідних спостережень, походження яких можна перевірити.

Розмежування контрольованого й операційного сценаріїв визначає допустиму силу висновку. У контрольованому сценарії відомі обидва об’єкти, правило підготовки та параметри внесення, тому можна перевіряти відмінність $D$ і збереження $S$ відносно заданого навантаження. Для невідомого входу така пара порівняння відсутня. Виявлена ознака в цьому разі характеризує властивість доступного представлення, але не відновлює автоматично ні $P_A$, ні $E_A$, ні $\theta_A$. Отже, формула~\eqref{eq:attack_embedding} задає модельний зв’язок між сутностями сценарію, а не процедуру визначення способу підготовки невідомого об’єкта.

У межах одного проходження кортеж $(g_j,\theta_j,\varphi_j,t_j)$ фіксує фактичну операцію над об’єктом. Одержані похідні дані не входять до складу $G$; їх пов’язують із відповідним елементом послідовності через відомості про походження. Наявність запису про виклик компонента не означає, що потрібне представлення сформовано повністю. Так само успішне формування представлення не підтверджує відсутності прихованого навантаження. Для подальшого причинного зіставлення потрібні відтворюваний опис операції та окремо збережений результат із явно встановленим статусом якості.

Повторна поява того самого компонента в $G$ не скорочує послідовність до одного етапу. Різні виклики можуть відбуватися за відмінних станів, параметрів і часових умов, а тому породжувати різні спостереження. Ідентифікатор запуску відокремлює їх від подій паралельних перевірок. Завдяки цьому ознаку пов’язують не лише з назвою компонента, а з конкретним виконанням операції над визначеним входом. Така деталізація потрібна для встановлення межі, у якій причинна належність події залишається перевірюваною.

Підготовлений об’єкт послідовно опрацьовують функціонально відмінні компоненти вебсистеми. Їх подано скінченною множиною $R_{web}=\{r_1,r_2,\ldots,r_{n_R}\}$, до якої можуть входити засоби приймання, первинної перевірки, структурного розбору, розкодування, перетворення, збереження, відображення та контрольованого спостереження. Фактичний порядок проходження визначено послідовністю $G=\langle(g_j,\theta_j,\varphi_j,t_j)\rangle_{j=1}^{n_G}$, де $g_j\in R_{web}$, $\theta_j$ описує стан компонента, $\varphi_j$ позначає виконану операцію, а $t_j$ є часом спостереження. Той самий компонент може виконувати кілька операцій і тому траплятися у послідовності неодноразово.

Основні умови доставки задають три функції:
\[
\begin{aligned}
V_f &\colon \mathrm B^{*}\rightarrow\{0,1\},\\
D   &\colon \mathrm B^{*}\times\mathrm B^{*}\rightarrow[0,\infty),\\
S   &\colon \mathrm B^{*}\times\mathrm G\rightarrow[0,1].
\end{aligned}
\]
де $\mathrm B=\{0,\ldots,255\}$, а $\mathrm G$ позначає множину допустимих послідовностей оброблення. Порогові значення задовольняють умови $\varepsilon\geq0$ і $\tau_A\in[0,1]$.

Кожна з трьох функцій описує окрему властивість сценарію. Значення $V_f$ стосується лише рішення конкретного правила приймання. Величина $D$ визначає відхилення між двома відомими байтовими послідовностями за обраним правилом порівняння. Величина $S$ характеризує збереження заданого навантаження після визначеної послідовності оброблення. Через різні області визначення ці значення не є взаємозамінними й не утворюють спільної числової шкали. Їх поєднання виражає одночасне виконання різних умов, а не сумарну оцінку загрози.

Можливість прихованої доставки до точки спостереження визначають дві взаємопов’язані умови.

Першою умовою є виконання правила приймання. Рівність $V_f(X_{mm}^{\prime})=1$ означає, що об’єкт допущено до подальшого оброблення. Тут $V_f$ позначає бінарну функцію перевірки формату. Ця умова підтверджує лише відповідність установленому правилу приймання і не встановлює безпечності вмісту.

Друга умова характеризує збереження прихованості та навантаження. Відмінність між початковим і підготовленим об’єктами обмежено умовою $D(X_{mm},X_{mm}^{\prime})\leq\varepsilon$. Збереження значущої частини навантаження під час проходження послідовності $G$ визначено нерівністю $S(P_A,G)\geq\tau_A$. Тут $D$ позначає функцію відхилення між об’єктами, $\varepsilon$ визначає допустимий рівень зміни, $S$ є оцінкою збереження навантаження, а $\tau_A$ задає порогове значення цієї оцінки. Виконання обох нерівностей характеризує можливість доставки, проте не доводить прояву прихованого впливу.

Поєднання результатів цих перевірок утворює кілька змістово різних станів. Якщо $V_f(X_{mm}^{\prime})=0$, штатний маршрут переривається на розглянутому правилі приймання незалежно від значень $D$ і $S$. Такий результат не спростовує наявності прихованого навантаження, а лише встановлює, що зафіксовану точку спостереження цим маршрутом не досягнуто. Якщо $V_f(X_{mm}^{\prime})=1$, але не виконано хоча б одну з двох порогових умов, проходження об’єкта саме по собі не підтверджує збереження заданого контрольованого навантаження. Лише одночасне виконання правила приймання й обох порогових умов установлює можливість його доставки в межах заданого сценарію.

Значення порогів $\varepsilon$ і $\tau_A$ належать опису конкретної контрольованої конфігурації. Зміна правила порівняння в $D$, складу послідовності $G$ або способу визначення $S$ змінює умови, за яких інтерпретують нерівності. Тому однакові числові значення, одержані за різних конфігурацій, не вважають тотожними доказовими підставами. Для повторення перевірки разом із результатом зберігають визначення функцій, пороги та опис послідовності оброблення.

Якщо в операційному сценарії немає еталонного $X_{mm}$ або відомого $P_A$, значення $D$ чи $S$ не відновлюють за непрямими ознаками. Наявні структурні, спектральні, зорові та подієві відхилення переходять до багаторівневого представлення з власними координатами й відомостями про походження. Це зберігає відмінність між перевіркою контрольовано підготовленого носія та аналізом невідомого входу й не перетворює припущення про спосіб підготовки на встановлений факт.

Доказове встановлення прихованого впливу потребує виконання двох окремих умов.

Першою умовою є відтворюваність послідовності оброблення. Належність спостереження до певної фази підтверджують записом відповідного елемента $g_j$, стану $\theta_j$, виконаної операції та часу. За відсутності цих відомостей не можна встановити, на якому етапі виникла або зникла досліджувана ознака.

Другою умовою є підтвердження причинної належності події. Подію враховують лише тоді, коли встановлено її зв’язок із досліджуваним об’єктом, конкретним запуском і відповідною фазою послідовності $G$. Часовий збіг без такого зв’язку не є достатньою підставою для висновку про прояв прихованого впливу.

Наведені умови розмежовують наявність прихованого навантаження, його збереження під час оброблення та спостережуваний прояв впливу. Структурне відхилення характеризує властивість контейнера. Якщо його спостережено після проходження певного компонента, цю властивість фіксують на відповідному етапі. Лише фактично спостережена причинно пов’язана реакція дає підстави встановити прихований вплив у межах доступних спостережень. Жоден із зазначених фактів окремо не встановлює наміру порушника або програмного змісту виявленого фрагмента.

Якщо відомості про частину послідовності $G$ відсутні, невідому ділянку фіксують як межу причинного висновку. Її не відновлюють за припущеннями. Це запобігає помилковому віднесенню сторонньої події до мультимедійного об’єкта лише за часовою близькістю.

Для тлумачення відсутності подій оцінюють не їх кількість, а покриття фаз, обов’язкових у чинній конфігурації контрольованого спостереження. Ця оцінка належить до супровідних відомостей і не змінює визначення $G$. Порожній журнал за підтвердженого покриття означає, що в межах спостережуваних фаз причинно пов’язаних подій не зареєстровано. Недоступний журнал або пропущена обов’язкова фаза не дають такого висновку. У першому випадку одержано спостережений порожній результат, а у другому спостереження відсутнє. Цю відмінність надалі зберігають за допомогою $X_b$ і складника $m_b$.

Послідовність від підготовки носія до одержання доступних спостережень узагальнено на рис.~\ref{fig:attack_model}. Разом з ознакою необхідно зберігати етап оброблення, на якому її виявлено, та відомості, що підтверджують її походження.

\begin{figure}[htbp]
\centering
\begin{tikzpicture}[x=1cm,y=1cm]
\node[dissertation block,text width=2.75cm] (object) at (-5.65,1.35)
  {Початковий об’єкт\\$X_{mm}$};
\node[dissertation block,text width=2.75cm] (payload) at (-5.65,-0.65)
  {Приховане навантаження\\$P_A$};
\node[dissertation block,text width=3.00cm] (prepare) at (-2.10,0.35)
  {Підготовка об’єкта\\$E_A(\Theta_A)$};
\node[dissertation block,text width=2.75cm] (prepared) at (1.20,0.35)
  {Підготовлений об’єкт\\$X_{mm}^{\prime}$};
\node[dissertation block,text width=3.15cm] (conditions) at (4.65,0.35)
  {Умови прихованості\\$V_f=1$, $D\leq\varepsilon$};
\node[dissertation block,text width=3.30cm] (route) at (4.65,-1.75)
  {Маршрут $G$\\$r_j\in R_{web}$};
\node[dissertation block,text width=4.40cm] (observations) at (0.55,-3.20)
  {Статичні та подієві спостереження\\$X_s,X_v,X_f,X_b$};

\coordinate (inputbus) at (-4.05,0.35);
\draw[thick] (object.east) -- (object.east -| inputbus) -- (inputbus);
\draw[thick] (payload.east) -- (payload.east -| inputbus) -- (inputbus);
\fill[black] (inputbus) circle (1.3pt);
\draw[dissertation arrow] (inputbus) -- (prepare.west);
\draw[dissertation arrow] (prepare.east) -- (prepared.west);
\draw[dissertation arrow] (prepared.east) -- (conditions.west);
\draw[dissertation arrow] (conditions.south) -- (route.north);
\draw[dissertation arrow] (route.south) |- (observations.east);
\end{tikzpicture}
\caption{Причинна схема формування спостережень за умов прихованої доставки мультимедійного об’єкта}
\label{fig:attack_model}
\end{figure}


Умови доставки описують можливість проходження та збереження навантаження, а умови доказового встановлення впливу визначають межу причинного висновку. Жодна з цих умов не доводить виконання шкідливого коду. Подальша перевірка потребує представлень, у яких збережено початковий носій, результати його контрольованого тлумачення та межі доступного спостереження.

\section{Формалізація мультимедійного об’єкта як багаторівневого інформаційного контейнера}
\label{sec:multilevel_container}

Мультимедійний об’єкт опрацьовують за кількома правилами, кожне з яких утворює окреме похідне представлення. Після розкодування одержують дані для зорового й часового аналізу, проте під час цієї операції службові ділянки можуть не відтворитися, а структурні помилки можуть бути виправлені засобом розкодування. Побайтовий розбір відтворює будову контейнера, але не розкриває зорового змісту. Журнал подій містить відомості про реакцію середовища, однак не підтверджує причинної належності кожного запису досліджуваному об’єкту. Тому одиницею аналізу є одна перевірка з підтвердженим походженням усіх використаних даних.

Ідентифікатор перевірки пов’язує початковий об’єкт, похідні дані, ознаки, події та аналітичний результат. Тотожність входу підтверджують незмінністю байтової послідовності. Належність похідних даних установлюють за збереженим правилом їх одержання. Єдиним безпосередньо зафіксованим джерелом є початкова послідовність байтів; довіра стосується лише її тотожності отриманому входу, а не змісту. Решту даних розглядають як результати спостереження із зазначенням застосованих засобів та їхніх версій.

Контракт кожного похідного представлення охоплює чотири складники: ідентифікатор початкового входу, відтворюване правило перетворення, тип одержаного результату та статус його якості. Відсутність хоча б одного складника обмежує подальше використання представлення. Без ідентифікатора неможливо підтвердити належність до поточної перевірки. Без правила й версії засобу не можна відтворити перетворення. Без типу результату числові значення не мають однозначного тлумачення. Без статусу якості повний результат неможливо відрізнити від часткового або технічно пошкодженого.

Якість представлення визначає межу його допустимого тлумачення, але не змінює належності до входу. Частково розкодований кадр може походити від потрібного $X_{mm}$, проте не охоплювати всю просторову область. Неповна структурна карта може містити достовірні координати розпізнаних елементів, але не давати підстав для висновку про нерозібрану частину контейнера. Тому показник покриття застосовують до фактично одержаної області, а невідомі ділянки не заповнюють припущеними значеннями.

Повний вхід моделі формують у п’ять послідовних кроків.

\begin{enumerate}
\item \textit{Збереження початкового носія.} Об’єкт $X_{mm}=(x_1,x_2,\ldots,x_n)$ зберігають як послідовність $n$ байтів зі значеннями з множини $\{0,\ldots,255\}$. Очищення, перекодування або автоматичне виправлення утворює похідну копію і не заміщує $X_{mm}$. Це правило є обов’язковим, оскільки перетворення може знищити досліджуване відхилення.

\item \textit{Одержання зорових і часових даних.} Контрольоване розкодування утворює $X_v=E_v(X_{mm},\theta_{\mathrm{dec}})$, де $E_v$ позначає оператор розкодування, а $\theta_{\mathrm{dec}}$ містить його параметри та версію. Для зображення $X_v$ складається з одного кадру, а для відео містить упорядковану послідовність кадрів. Частотне представлення $X_f$ і часове представлення $X_t$ формують під час формування ознак. Вони є похідними від розкодованих даних і не замінюють $X_v$.

\item \textit{Відтворення будови контейнера.} Структурне представлення $X_s=\langle s_1,s_2,\ldots,s_{n_s}\rangle$ зберігає порядок і вкладеність виявлених елементів. Запис $s_k=\langle c_k,a_k,b_k,w_k,v_k\rangle$ містить тип елемента $c_k$, напіввідкритий байтовий проміжок $[a_k,b_k)$, посилання $w_k$ на батьківський елемент і результат $v_k$ перевірки форматних обмежень. Права межа $b_k$ до проміжку не входить, тому його довжина дорівнює $b_k-a_k$. Пов’язані службові значення утворюють $X_m$. Байтові координати пов’язують кожне структурне відхилення з відповідною ділянкою початкового носія.

\item \textit{Відбір причинно пов’язаних подій.} Представлення ознак подій і поведінки $X_b=\langle e_1,e_2,\ldots,e_T\rangle$ містить лише події, належність яких до поточної перевірки підтверджено. Для події $e_t$ зберігають дію, ініціатора, ресурс, результат і час. Супровідний запис містить фазу оброблення, ідентифікатор запуску та стан середовища. Часовий збіг без причинного зв’язку не є підставою для включення події до $X_b$.

\item \textit{Фіксація доступності та якості.} Для трьох груп даних зберігають маску доступності $m_A=(m_s,m_v,m_b)\in\{0,1\}^3$. Її складники відповідають структурним даним, результатам спектрального та зорового аналізу, а також даним про події й поведінку. Нульове значення означає відсутність придатного спостереження, а не підтвердження норми. Причину недоступності, повноту спостереження та версії засобів уносять до супровідних відомостей $C$.
\end{enumerate}

Маска $m_A$ і відомості $C$ виконують взаємодоповнювальні функції. Маска вказує, чи можна використовувати групу даних у формальному обчисленні. Відомості $C$ пояснюють, чому група доступна або недоступна, якою є її повнота та за яких умов її одержано. Однакова маска двох перевірок не означає однакової якості спостережень. Наприклад, значення $m_v=1$ може супроводжувати як повністю розкодовану послідовність кадрів, так і допустимо неповне подання з указаним покриттям. Різницю між цими випадками зберігають у $C$ і враховують під час формування упевненості.

Доступність установлюють після перевірки походження й мінімальної придатності результату. Сам факт створення проміжного об’єкта не є достатньою підставою для значення маски, що дорівнює одиниці. Якщо засіб завершив роботу, але сформований результат не пов’язано з поточним входом або його тип не відповідає очікуваному контракту, відповідну групу не допускають до обчислення. Технічний артефакт при цьому можна зберегти для діагностики, проте він не стає складником $X_A$.

Спостережений нуль, порожній результат і недоступність належать до різних станів. Нульове числове значення ознаки є результатом вимірювання в доступному представленні. Порожня послідовність $X_b=\varnothing$ є результатом повного спостереження обов’язкових фаз без відібраних причинно пов’язаних подій. Недоступність означає, що вимірювання або спостереження не забезпечує мінімального контракту. Змішування цих станів змінило б не лише числове значення, а й логіку висновку, тому їх розрізняють до завершення аналізу.

Координати забезпечують простежуваність усередині окремого представлення. Байтовий проміжок структурного елемента безпосередньо посилається на $X_{mm}$. Просторова або часова область посилається на $X_v$ чи $X_t$ і лише через збережене правило розкодування пов’язується з початковими байтами. Подієва координата посилається на фазу й час у $G$. Ці системи координат не ототожнюють. Їх узгоджують настільки, наскільки це дозволяє відтворюване перетворення, а невизначеність відображення зберігають як межу локалізації.

Суперечність між представленнями не усувають вибором одного з них без формального правила. Структурна карта може вказувати на відхилення за успішного розкодування, оскільки різні оператори перевіряють різні властивості. Так само невдале розкодування не скасовує доступних байтових координат. За таких умов обидва результати зберігають разом із власним походженням, а характер суперечності передають до ознакового ядра. Це дає змогу відокремити різні спостережувані властивості від технічної помилки узгодження.

Окремі представлення придатні до спільного аналізу лише тоді, коли їх віднесено до одного входу й відомі умови їх одержання. Статичні дані згруповано як $X_c=\langle X_{mm},X_v,X_s,X_m\rangle$. Повний вхід моделі визначено формулою~\eqref{eq:prepared_packet}:
\begin{equation}
X_A=\langle X_c,X_b,m_A,C\rangle,
\label{eq:prepared_packet}
\end{equation}

де $X_A$ позначає повний вхід моделі; $X_c$ є сукупністю статичних даних; $X_b$ містить причинно пов’язані події; $m_A$ позначає маску доступності; $C$ містить відомості про походження, параметри та якість даних.

Формула~\eqref{eq:prepared_packet} пов’язує спостереження з однією перевіркою, зберігаючи відмінність між статичними даними та даними про події й поведінку. Статичні ознаки формують без урахування реакції середовища. Ознаки подій і поведінки приєднують після підтвердження причинної належності подій. Отже, склад $X_A$ визначає вимоги до відтворюваності аналізу, не встановлюючи способу програмної реалізації.

Граничні випадки мають явні позначення. Якщо журнал повністю охоплює обов’язкові фази, але причинно пов’язаних подій не виявлено, встановлюють $X_b=\varnothing$ і $m_b=1$; порожня послідовність у цьому разі є спостереженим результатом. Значення $m_b=0$ використовують лише за недоступного журналу, недостатнього фазового покриття або непідтвердженого походження записів. Якщо розкодування завершилося без результату, структурні дані можуть залишатися доступними за $m_v=0$. Похідне представлення з непідтвердженим походженням не включають до $X_A$, навіть якщо його значення можна прочитати. Доступне, але неповне представлення супроводжують показником покриття та зазначенням меж його тлумачення.

Контракт $X_A$ допускає неповноту, але не допускає прихованої невизначеності. Неповнота є допустимою, коли відомо, якої групи даних бракує, з якої причини та які інші групи залишаються придатними. Невизначеність стає недопустимою, коли неможливо встановити належність даних до входу, версію перетворення або значення маски. У першому випадку подальший аналіз може сформувати явно частковий результат. У другому інтеграцію припиняють, оскільки немає відтворюваного способу визначити склад вхідних підстав.

Оновлення засобу розкодування або структурного розбору утворює нові похідні спостереження навіть за незмінного $X_{mm}$. Попередні результати не перезаписують, оскільки відмінність може бути спричинена версією оператора, а не зміною об’єкта. Порівняння таких запусків виконують через спільний ідентифікатор початкових байтів та окремі записи параметрів. Завдяки цьому зміна інструментального тлумачення не маскується під зміну властивостей носія.

Формування статичних складників і даних про події та поведінку у складі повного входу показано на рис.~\ref{fig:multilevel_container}. Стрілки від $X_{mm}$ до $X_v$, $X_s$ і $X_m$ позначають контрольоване одержання похідних представлень, а шлях до $X_b$ відображає відбір подій, зареєстрованих під час оброблення того самого об’єкта. До ознакового аналізу ці дані передають разом із маскою доступності та відомостями про походження і якість.

\begin{figure}[htbp]
\centering
\begin{tikzpicture}[x=1cm,y=1cm]
\node[dissertation block,text width=3.05cm] (object) at (-5.70,1.20)
  {Незмінний об’єкт\\$X_{mm}$};
\node[dissertation block,text width=3.35cm] (visual) at (-2.05,2.25)
  {Зорові дані\\$X_v$};
\node[dissertation block,text width=3.35cm] (structural) at (-2.05,0.15)
  {Структурні дані\\$X_s,X_m$};
\node[dissertation block,text width=4.15cm] (static) at (1.70,1.20)
  {Статичні дані після попереднього опрацювання\\$X_c=\langle X_{mm},X_v,X_s,X_m\rangle$};
\node[dissertation block,text width=3.40cm] (events) at (-2.05,-2.20)
  {Причинно пов’язані події\\$X_b$};
\node[dissertation block,text width=3.20cm] (context) at (1.70,-1.35)
  {Доступність і супровідні\\відомості $(m_A,C)$};
\node[dissertation block,text width=3.20cm] (input) at (5.45,0.15)
  {Повний вхід моделі\\$X_A$};

\draw[dissertation arrow] (object.east) -- (visual.west);
\draw[dissertation arrow] (object.east) -- (structural.west);
\draw[dissertation arrow] (object.north) -- ++(0,1.70) -| (static.north);
\draw[dissertation arrow] (visual.east) -- (static.west);
\draw[dissertation arrow] (structural.east) -- (static.west);
\draw[dissertation arrow] (static.east) -- (input.west);
\draw[dissertation arrow] (events.east) -- ++(0.55,0) -| ([xshift=-0.45cm]input.south);
\draw[dissertation arrow] (context.east) -- (input.west);
\end{tikzpicture}
\caption{Походження статичних і подієво-поведінкових даних повного входу моделі}
\label{fig:multilevel_container}
\end{figure}


Багаторівневе представлення поєднує незмінний початковий об’єкт із похідними спостереженнями, для яких відомі координати, походження та повнота. Для узгодженого аналізу ці дані перетворюють на порівнювані ознаки зі збереженням зв’язку з джерелом.

\section{Ознакове ядро моделі мультимодального представлення}
\label{sec:feature_space}

Багаторівневі дані мають різну фізичну природу та різні системи координат. Їх передчасне об’єднання може спричинити числове переважання однієї групи ознак і втрату зв’язку між оцінкою та ділянкою об’єкта, з якої цю оцінку одержано. Тому ознакове ядро приводить дані до порівнюваного вигляду, зберігаючи відмінність між групами до етапу їх обґрунтованого узгодження.

Порівнюваність ознак потребує зафіксованої схеми для кожної групи. Схема визначає порядок складників, їхню розмірність, правило обчислення, допустимі області значень і параметри попереднього перетворення. Однаковий розмір двох векторів не підтверджує їхньої сумісності, якщо відрізняються версії схеми, спосіб нормування або правило відбору ділянок. Тому версія ознак є частиною походження нарівні з ідентифікатором входу й версією засобу одержання представлення.

Числова форма не скасовує предметного змісту ознаки. Складник $Z_s$ характеризує структурні властивості й має байтову локалізацію, складник $Z_v$ характеризує спектральні та зорові властивості й має просторово-часову локалізацію, а $Z_b$ описує причинно допущені події. Їх можна подати дійсними числами, але одиниця, нуль або однакова відстань між векторами мають різне предметне тлумачення для кожної групи. Узгодження виконується лише після збереження цих відмінностей у структурі результату.

Перетворення даних на ознаки організовано за чотирма послідовними кроками.

Формування структурних ознак починається з опрацювання карти $X_s$, яка має змінну довжину й залежить від формату, тому безпосереднє порівняння об’єктів за нею неможливе. Функція $F_s$ перетворює $X_s$ і службові значення $X_m$ на числове представлення $Z_s=F_s(X_s,X_m;\theta_s)$. Функція $G_s$ утворює множину байтових координат $L_s=G_s(X_s;\theta_s)$. Представлення $Z_s$ використовують для зіставлення об’єктів, а $L_s$ застосовують для визначення ділянок початкового носія, які обґрунтовують одержану оцінку.

Спектральні й зорові ознаки формують з урахуванням того, що структурно коректний контейнер може містити нетипові закономірності в розкодованому вмісті. Тому зорове представлення доповнюють частотним $X_f$ і, для відео, часовим представленням $X_t$. Функція $F_v$ формує $Z_v=F_v(X_v,X_f,X_t;\theta_v)$. Функція $G_v$ визначає просторово-часові координати за покадровими ознаками, оцінками, показниками технічної якості та часовими мітками: $L_v=G_v(\{Z_{v,j},\rho_{v,j},q_j,t_j\}_{j=1}^{n_v},\tau_v;\Theta_H)$, де $\tau_v$ позначає поріг відбору ділянок із локальними відхиленнями. Цей поріг належить до параметрів першого методу й конкретизується в розділі~\ref{chap:methods_technology}. Для зображення встановлюють $X_t=\varnothing$ як коректне незастосовне значення, відмінне від недоступного спостереження.

Ознаки подій і поведінки формують із послідовності $X_b$, яка має змінну довжину й містить категорійні та часові значення. Після перевірки причинної належності подій функція $F_b$ формує представлення $Z_b=F_b(X_b;\theta_b)$. Відсутність події можна тлумачити лише тоді, коли відповідна фаза була доступною для спостереження. Тому $Z_b$ розглядають разом зі складником $m_b$ та показниками повноти журналу.

Завершальним кроком є узгодження походження. Байтові координати $L_s$ і просторово-часові координати $L_v$ зіставляють за збереженим правилом одержання розкодованого представлення. Відсутність прямої відповідності не спростовує результатів жодної складової аналізу, оскільки розкодування може мати нелокальний характер. У такому разі обмежують лише висновок про спільне джерело відхилень.

Кожна функція формування ознак повинна зберігати можливість повернення від числового складника до спостереження, яке його обґрунтовує. Для структурної групи таким зв’язком є елемент $X_s$ і його байтовий проміжок. Для спектральної та зорової групи це кадр, просторова область, часовий інтервал і параметри перетворення. Для подієво-поведінкової групи такими спостереженнями є відібрані записи, фаза $G$ та правило причинного допуску. Якщо числове значення не має такого посилання, його можна зберегти як технічний проміжний результат, але не використовувати як перевірювану підставу пояснення.

Узгодження походження не вимагає штучної точкової відповідності між усіма системами координат. Одна байтова область може впливати на значну просторову або часову частину розкодованого вмісту, а одна подія може стосуватися операції над цілим об’єктом. Тому $L_s$, $L_v$ і посилання на події зберігають власну гранулярність. Спільне джерело встановлюють лише тоді, коли відтворюване правило перетворення дає для цього достатні підстави. Інакше результати залишаються пов’язаними з одним входом, але не подаються як локалізований спільний прояв.

Для спільного оцінювання статичних ознак може бути утворено похідне представлення $Z_c$. Область визначення функції $F_C$ задають разом із маскою $m_c=(m_s,m_v)$, тому недоступне представлення не підміняють нульовим вектором. Для фіксованої розмірності $d_c$ об’єднання з урахуванням маски визначено формулою~\eqref{eq:masked_static_fusion}:
\begin{equation}
\begin{gathered}
F_C(Z_v,Z_s,(1,1);\theta_C)=F_C^{11}(Z_v,Z_s;\theta_C^{11})\in\mathrm R^{d_c},\\
F_C(Z_v,\bot,(0,1);\theta_C)=F_C^{01}(Z_v;\theta_C^{01})\in\mathrm R^{d_c},\\
F_C(\bot,Z_s,(1,0);\theta_C)=F_C^{10}(Z_s;\theta_C^{10})\in\mathrm R^{d_c},\\
F_C(\bot,\bot,(0,0);\theta_C)=\bot,
\end{gathered}
\label{eq:masked_static_fusion}
\end{equation}

де $\mathrm R$ позначає множину дійсних чисел; $\theta_C^{ab}$ є окремою конфігурацією із зафіксованою версією для маски $(a,b)$; $Z_c$ є числовою проєкцією, сформованою з урахуванням маски. Проєкцію за масок $(1,0)$ або $(0,1)$ обчислюють лише за наявності відповідної зафіксованої конфігурації; інакше встановлюють $Z_c=\bot$, зберігаючи доступні $Z_s$ або $Z_v$ у статичному описі. У підрозділі~\ref{sec:behavior_alignment} $Z_c$ застосовано в конфігурації методу, що ґрунтується на визначених правилах. Представлення $Z_c$ не замінює $Z_v$ і $Z_s$, оскільки однакове його значення може відповідати різним поєднанням ознак та різній повноті спостережень. Тому від моделі до подальшого опрацювання передають статичний опис, визначений формулою~\eqref{eq:static_profile}:

Відображення $F_C^{11}$, $F_C^{01}$ і $F_C^{10}$ є окремими маскозалежними контрактами, навіть якщо їхні результати належать одному простору $\mathrm R^{d_c}$. Для кожного контракту фіксують власні параметри й умови застосування. Конфігурацію повного складу не застосовують до неповного входу через заповнення відсутнього вектора нулями, середніми значеннями або значеннями іншого об’єкта. Таке заповнення змінило б предметний стан «спостереження відсутнє» на штучне числове спостереження.

Спільна розмірність $d_c$ забезпечує технічну сумісність виходу, але не встановлює рівносильності проєкцій. Значення, одержані за різних масок, можуть мати відмінний розподіл і відмінну доказову основу. Тому разом із $Z_c$ завжди зберігають маску й версію $\theta_C^{ab}$. Порівняння двох проєкцій допустиме лише в межах правила, яке прямо враховує цю відмінність. Без такого правила однакове числове значення не має спільного каліброваного тлумачення.

Недоступність однієї статичної групи не зменшує значень доступної групи й не спростовує її локальних результатів. Вона змінює склад підстав, за якими можна сформувати спільну проєкцію. Якщо маскозалежної конфігурації немає, модель зберігає $Z_s$ або $Z_v$ у $H$, але не створює $Z_c$. Отже, відмова від проєкції є контрольованим станом контракту, а не втратою вже одержаних ознак.
\begin{equation}
H=\langle Z_v,Z_s,\rho_v,\rho_s,L_v,L_s,m_c,C_H\rangle,
\label{eq:static_profile}
\end{equation}

де $H$ позначає статичний опис об’єкта; $Z_v$ і $Z_s$ містять відповідно сукупність спектральних і зорових ознак та сукупність структурних ознак; $\rho_v$ і $\rho_s$ є відповідними частковими оцінками; $L_v$ і $L_s$ визначають просторово-часові та байтові координати; $m_c=(m_s,m_v)$ позначає маску доступності статичних складових; $C_H$ містить відомості про умови формування статичного опису.

Спільне збереження числових значень, часткових оцінок, координат їхнього походження та маски доступності робить підстави результату простежуваними під час переходу до методу поведінкового узгодження. Представлення $X_b$ до $H$ не входить, оскільки додавання подій на цьому етапі змішало б властивості носія з реакцією середовища.

Статичний опис $H$ є межею між формуванням ознак носія та подальшим причинним зіставленням. До цієї межі часткові оцінки $\rho_v$ і $\rho_s$ характеризують лише відповідні властивості доступних статичних представлень. Вони не містять висновку про реакцію компонента й не підтверджують прихованого впливу. Завдяки цьому подальше приєднання $Z_b$ можна перевірити як окрему операцію, а не як неявне використання подій під час формування статичних ознак.

Для передавання $H$ потрібна внутрішня узгодженість його складників. Значення $m_c$ мають відповідати фактичній наявності $Z_s$, $Z_v$, $L_s$ і $L_v$, а $C_H$ має містити версії схем та ідентифікатор входу. Якщо часткова оцінка доступна без відповідного ознакового вектора або координат, статичний опис не відповідає визначеному контракту. Такий набір не доповнюють на наступному етапі, а повертають до перевірки формування ознак.

Після окремого формування трьох груп ознак їх можна об’єднувати лише за належності одному входу, сумісності умов одержання та підтвердженого походження. Відсутнє спостереження при цьому не прирівнюють до нульового значення, а числову подібність не тлумачать як причинний зв’язок. Залежність між статичним описом, ознаками подій і поведінки та повним представленням визначено формулою~\eqref{eq:integrated_representation}:
\begin{equation}
Z=F_I(H,Z_b,m_A,C_I;\theta_I)=\langle z,L,m_A,C_I\rangle,
\label{eq:integrated_representation}
\end{equation}

де $Z$ позначає інтегроване представлення; $F_I$ є функцією узгодженого об’єднання; $z$ є числовим вектором з окремо позначеними складниками трьох груп ознак та їх частковими оцінками; $L$ містить сукупність байтових і просторово-часових координат із посиланнями на події; $m_A=(m_s,m_v,m_b)$ позначає маску доступності структурної, спектрально-зорової та подієво-поведінкової груп; $C_I$ містить доповнені відомості про якість, походження та параметри інтеграції; $\theta_I$ є версованою конфігурацією інтеграції. У $C_I$ зберігають доступність конкретних похідних представлень, незастосовність $X_t$ і причини неповноти; ці відомості не утворюють додаткових складників маски $m_A$.

Перед обчисленням перевіряють, що статична проєкція маски $m_A|_{\{s,v\}}$ збігається з $m_c$, а статична частина $C_I$ збігається з $C_H$. За розбіжності ідентифікаторів, контрольних сум або версій інтеграцію припиняють. У такий спосіб повторне подання маски й відомостей до $F_I$ доповнює статичний опис поведінковими даними, але не допускає підміни вже сформованої частини $H$.

Формула~\eqref{eq:integrated_representation} не передбачає простого усереднення ознак. Узгоджене підтвердження, взаємне доповнення, суперечність і недостатність підстав мають залишатися різними станами. Тому однакові значення $z$, одержані за різної повноти або суперечливості даних, не вважають рівносильними. Відмінність між ними зберігають у $L$, $m_A$ і $C_I$.

Узгоджене підтвердження виникає тоді, коли доступні групи вказують на сумісний предметний стан і мають підтверджене походження. Взаємне доповнення означає, що різні групи описують різні властивості одного входу без прямої суперечності. Суперечність фіксують, коли доступні результати не можуть бути одночасно витлумачені в межах однієї конфігурації. Недостатність підстав установлюють за відсутності мінімального складу або якості спостережень. Ці стани описують відношення між підставами; вони не є додатковими класами загрози й не замінюють $\hat k$.

Функція $F_I$ не повинна перетворювати суперечність на середнє значення, яке приховує напрям локальних оцінок. Якщо одна група містить виразне відхилення, а інша доступна група його не підтримує, у $z$ зберігають обидва результати, а в $C_I$ фіксують ознаку їхньої неузгодженості. Подальше рішення може врахувати цю неузгодженість через $c_A$ та пояснення, але не змінює вихідних часткових оцінок. Так зберігається можливість перевірити, яка саме група визначила аналітичний результат.

Взаємне доповнення також не тлумачать як причинність. Структурна й зорова ознаки можуть стосуватися одного входу та різних його ділянок. Подієва ознака може виникнути під час оброблення цього входу, але причинний допуск визначає фазова належність, а не числова близькість. Тому $F_I$ поєднує типізовані відомості, зберігаючи вже встановлені причинні зв’язки, але не створює новий причинний зв’язок лише внаслідок об’єднання.

Для однозначного тлумачення області застосування основних функцій визначено так. Відображення $E_A:\mathrm{B}^{*}\times\mathrm{B}^{*}\times\mathrm T_A\rightarrow\mathrm{B}^{*}$ перетворює початкову й приховану байтові послідовності та конкретний кортеж параметрів $\theta_A\in\mathrm T_A$ на підготовлену байтову послідовність, де $\mathrm{B}=\{0,\ldots,255\}$.

Функції $F_s$ і $F_v$ перетворюють доступні структурні та зорово-частотні представлення на скінченновимірні дійсні вектори $Z_s\in\mathrm R^{d_s}$ і $Z_v\in\mathrm R^{d_v}$. Розмірності $d_s$ і $d_v$ визначає версія схеми ознак. Функція $F_b$ перетворює скінченну послідовність унормованих подій на $Z_b\in\mathrm R^{d_b}$.

Для недоступного числового представлення введено спеціальне ненумеричне значення $\bot$. Для $k\in\{s,v,b\}$ виконується $Z_k\in\mathrm R^{d_k}\cup\{\bot\}$, причому $m_k=0$ тоді й лише тоді, коли $Z_k=\bot$, а $m_k=1$ тоді й лише тоді, коли $Z_k\in\mathrm R^{d_k}$. Порожній, але повністю спостережений журнал має $m_b=1$ і відтворюване числове представлення $Z_b=F_b(\varnothing;\theta_b)$; це значення не є тотожним недоступності. Аналогічно координати недоступної складової позначають $\bot$, тоді як $X_t=\varnothing$ для зображення означає незастосовність часової складової.

Функція $F_I$ приймає $H$, $Z_b\in\mathrm R^{d_b}\cup\{\bot\}$, маску $m_A=(m_s,m_v,m_b)\in\{0,1\}^{3}$, відомості $C_I$ та конфігурацію $\theta_I$ і повертає $Z$. Якщо $m_b=0$, то $Z_b=\bot$, а неповне інтегроване представлення визначають формулою~\eqref{eq:partial_integrated_representation}:
\begin{equation}
Z^{\partial}=F_I\bigl(H,\bot,(m_s,m_v,0),C_I;\theta_I\bigr)
=\langle z_{sv},L_{sv},(m_s,m_v,0),C_I^{\partial}\rangle,
\label{eq:partial_integrated_representation}
\end{equation}

де $z_{sv}$ і $L_{sv}$ містять лише доступні статичні складові, а $C_I^{\partial}$ визначає статус неповноти та її документовану причину. Недоступні складові не замінюють нульовими спостереженнями. За $m_s=m_v=0$ функція $F_I$ не визначена й аналітичний результат не формують. Відомості $C_I$ містять ідентифікатор входу, версії засобів і схем, часові межі, показники якості та причини недоступності. Розмірності та порядок складників фіксують в окремому описі параметрів. За відсутності такого опису наведені залежності визначають лише формалізовані вимоги до представлень і не є достатнім описом конкретного запуску.

Часткове представлення $Z^{\partial}$ є повноправним результатом формування за явно неповного складу, але не еквівалентом повного $Z$. Його предметний зміст обмежено статичними складовими, а відсутність $Z_b$ відображено одночасно в масці й у $C_I^{\partial}$. Такий результат можна передати до функції, контракт якої допускає відповідну маску, проте не можна тлумачити як спостережену відсутність подій. Межа між двома випадками зберігається до аналітичного стану.

Мінімальною передумовою інтеграції є доступність принаймні однієї статичної групи. Це обмеження випливає з предмета моделі: аналітичний результат має стосуватися властивостей конкретного мультимедійного носія. Події без доступного статичного опису не утворюють інтегрованого представлення, навіть якщо належать певному запуску, оскільки немає мінімальної ознакової основи для стану об’єкта. Водночас доступність однієї статичної групи не робить результат повним; її склад і межі прямо фіксують у масці.

Відмова інтеграції може мати різні причини, які не об’єднують в один невизначений стан. Розбіжність ідентифікаторів означає, що складники належать різним входам або запускам. Несумісність схем означає, що числові представлення не мають спільного контракту. Непідтверджене походження не дає змоги перевірити зв’язок із входом. Непідтримувана маска вказує на відсутність відповідної конфігурації. Недостатня якість означає, що формально одержаний результат не досягає мінімальної придатності. Точна причина відмови визначає, чи потрібне повторне формування окремого представлення, узгодження версій або новий запуск усієї перевірки.

Подібність ознак сама по собі не підтверджує причинного зв’язку. Для причинного тлумачення потрібні спільний початковий об’єкт, сумісна послідовність оброблення та підтверджена належність подій. Якщо хоча б одну з цих умов не виконано, ознака залишається описовою і не посилює висновку про прояв прихованого впливу.

Вхідні дані, перевірку їх узгодженості та формування повного $Z$ або часткового $Z^{\partial}$ подано на рис.~\ref{fig:feature_space}. За розбіжності ідентифікаторів, контрольних сум, версій схем, походження подій або непідтримуваної маски інтеграцію припиняють. Підтримувану неповноту позначають маскою і не компенсують значеннями інших груп ознак.

\begin{figure}[H]
\centering
\begin{tikzpicture}[x=1cm,y=1cm,scale=0.88,transform shape,
  decision/.style={diamond,draw=black,line width=0.8pt,fill=white,
    aspect=1.35,align=center,font=\footnotesize,inner sep=1pt}]
\node[dissertation small,text width=3.10cm] (static) at (-6.50,2.20)
  {Статичний опис\\[-1pt]
   $H=\langle Z_v,Z_s,$\\
   $\rho_v,\rho_s,L_v,L_s,$\\
   $m_c,C_H\rangle$};
\node[dissertation small,text width=2.70cm] (behavior) at (-6.50,0.50)
  {Ознаки подій\\[-1pt]
   і поведінки $Z_b$};
\node[dissertation small,text width=3.10cm] (context) at (-6.50,-1.15)
  {Маска доступності,\\[-1pt]
   відомості й конфігурація\\
   $m_A,C_I,\theta_I$};

\node[dissertation small,text width=2.40cm,minimum height=19mm] (validation) at (-3.45,0.50)
  {Перевірка контракту\\
   ідентифікатор;\\
   контрольна сума;\\
   версії схем; походження;\\
   причинна\\
   належність;\\
   маска; якість};
\node[decision,minimum width=1.55cm,minimum height=9mm] (condition) at (-0.55,0.50)
  {Вимоги\\виконано?};
\node[dissertation small,text width=2.55cm,minimum height=11mm] (stop) at (-0.55,-1.75)
  {Інтеграцію припинено;\\
   причину відмови зафіксовано};
\node[dissertation small,text width=1.75cm,minimum height=10mm] (integration) at (2.10,0.50)
  {Узгоджене\\об’єднання\\$F_I$};
\node[decision,minimum width=1.45cm,minimum height=9mm] (composition) at (4.50,0.50)
  {Вхід\\повний?};
\node[dissertation small,text width=3.00cm,minimum height=11mm] (full) at (7.85,1.65)
  {Повне\\
   представлення\\[-1pt]
   $Z=\langle z,L,m_A,C_I\rangle$};
\node[dissertation small,text width=3.60cm,minimum height=13mm] (partial) at (7.85,-1.20)
  {Підтримуване часткове представлення\\[-1pt]
   $Z^{\partial}=\langle z_{sv},L_{sv},$\\
   $(m_s,m_v,0),C_I^{\partial}\rangle$};

\coordinate (inputbus) at (-4.80,0.50);
\draw[thick] (static.east) -- (static.east -| inputbus);
\draw[thick] (behavior.east) -- (behavior.east -| inputbus);
\draw[thick] (context.east) -- (context.east -| inputbus);
\draw[thick] (static.east -| inputbus) -- (context.east -| inputbus);
\fill[black] (inputbus) circle (1.3pt);
\draw[dissertation arrow] (inputbus) -- (validation.west);
\draw[dissertation arrow] (validation.east) -- (condition.west);
\draw[dissertation arrow] (condition.east) -- node[midway,above,font=\footnotesize,fill=white,inner sep=0.5pt]{так} (integration.west);
\draw[dissertation arrow] (condition.south) -- node[midway,right,font=\footnotesize,fill=white,inner sep=0.5pt]{ні} (stop.north);
\draw[dissertation arrow] (integration.east) -- (composition.west);
\coordinate (resultbus) at (5.75,0.50);
\draw[thick] (composition.east) -- (resultbus);
\draw[dissertation arrow] (resultbus) |- node[pos=0.20,right,font=\footnotesize,fill=white,inner sep=0.5pt]{так} (full.west);
\draw[dissertation arrow] (resultbus) |- node[pos=0.20,right,font=\footnotesize,fill=white,inner sep=0.5pt]{ні} (partial.west);
\end{tikzpicture}
\caption{Перевірка контракту та формування повного або часткового інтегрованого представлення}
\label{fig:feature_space}
\end{figure}


Ознакове ядро утворює простір узгодження, у якому числові значення зберігають зв’язок із координатами, доступністю та походженням. Підтверджувальне свідчення, суперечливе свідчення й відсутнє спостереження за таких умов мають різне тлумачення під час прийняття рішення.

\section{Критерії формування рішення щодо ознак аномального вмісту}
\label{sec:decision_criteria}

Інтегроване представлення може містити виразні, але неповні свідчення або помірні, проте узгоджені спостереження. Ці випадки відрізняються надійністю підстав і не можуть бути вичерпно описані одним числом. Тому аналітичний результат окремо містить визначений клас стану, інтегральну оцінку доказових підстав, упевненість і пояснення.

Складники аналітичного результату відповідають на різні питання. Клас $\hat k$ називає стан, установлений за правилами конфігурації. Значення $\rho_A$ характеризує виразність доступних підстав відносно порога. Упевненість $c_A$ відображає, наскільки повними, якісними й несуперечливими є ці підстави. Пояснення $\Pi$ дає змогу простежити результат до ознак, координат, маски та параметрів. Зміна одного складника не повинна неявно підміняти інший, оскільки висока виразність може поєднуватися з неповним спостереженням, а узгоджені слабкі підстави можуть не досягати порога.

Відокремлення цих складників визначає межу автоматичного тлумачення. Значення $\hat y$ придатне для передавання до зовнішнього правила реагування, але не містить повного обґрунтування. Клас без $\rho_A$ не показує відстані від порога, а $\rho_A$ без $c_A$ не характеризує надійності підстав. Пояснення без версії $\mathrm T_D$ не забезпечує відтворюваності обчислення. Тому мінімальним результатом є цілісний кортеж $S_A$, а не окреме вибране число.

Клас характеризує встановлений тип стану. Інтегральна оцінка доказових підстав відображає виразність сукупності доступних свідчень щодо можливого прихованого впливу, а упевненість відображає їх повноту, якість і несуперечливість. Пояснення пов’язує результат із ділянками об’єкта та умовами перевірки. Захисна дія до аналітичного стану не входить, оскільки її визначає інформаційна технологія з урахуванням чинного правила реагування.

Нехай $\mathrm K=\{\mathrm K_0,\mathrm K_1,\ldots,\mathrm K_{n_A},\mathrm K_U\}$ позначає множину класів стану. Клас $\mathrm K_0$ відповідає недосягненню встановленого порога інтегральної оцінки доказових підстав, класи $\mathrm K_1,\ldots,\mathrm K_{n_A}$ відповідають визначеним типам прихованої загрози, а $\mathrm K_U$ відповідає стану, за якого наявні підстави для захисної дії, але конкретний тип загрози не встановлено. Параметри прийняття рішення утворюють множину $\mathrm T_D$, яка містить поріг інтегральної оцінки $\tau$, поріг класової оцінки $\tau_p$, мінімальний розрив $\delta_p$, базові вагові коефіцієнти, порядок розв’язання однакових класових оцінок і версії правил обчислення упевненості та пояснення. Перехід від інтегрованого представлення до аналітичного стану визначено формулою~\eqref{eq:analytic_state}:
\begin{equation}
S_A=F_D(Z,\mathrm T_D)=\langle p,\hat y,\hat k,\rho_A,c_A,\Pi\rangle,
\label{eq:analytic_state}
\end{equation}

де $S_A$ позначає аналітичний стан мультимедійного об’єкта; $F_D$ є функцією прийняття рішення; $Z$ позначає інтегроване представлення; $\mathrm T_D$ є множиною параметрів прийняття рішення; $p=(p_1,\ldots,p_{n_A})\in\mathrm P^{n_A-1}$ є нормованим вектором оцінок належності до класів можливого впливу, для якого $p_j\geq0$ і $\sum_{j=1}^{n_A}p_j=1$; $\hat y\in\{0,1\}$ позначає двійковий висновок; $\hat k\in\mathrm K$ є визначеним класом; $\rho_A$ є інтегральною оцінкою доказових підстав щодо можливого прихованого впливу, обчисленою для конкретної конфігурації ознак, маски доступності й параметрів рішення; $c_A$ позначає упевненість; $\Pi$ є перевірюваним поясненням.

Оцінки $p_j$ характеризують відносну відповідність ознак визначеним класам і не є статистичними ймовірностями без окремої перевірки калібрування. Величина $\rho_A$ також не є статистичною ймовірністю, очікуваним збитком або універсально порівнюваною мірою. Вона призначена для порівняння з порогом, установленим для конкретної конфігурації.

Нормування вектора $p$ установлює лише відносний розподіл класових оцінок у межах одного запуску. Воно не перетворює ці оцінки на частоти подій і не дає підстав зіставляти їх між різними версіями схеми без окремої перевірки. Високе $p_j$ означає переважання відповідного класу серед визначених альтернатив, але конкретний клас призначають лише разом з умовами порога й мінімального розриву. Отже, нормування забезпечує однозначність формального вибору, а не статистичну інтерпретацію.

Інтегральна оцінка $\rho_A$ має іншу функцію: вона узгоджує доступні часткові оцінки для порівняння з $\tau$. Її значення залежить від складу $I_{\rho}$, базових ваг і маски. Тому однакове $\rho_A$ може бути одержано з різних комбінацій підстав. Відмінність між ними не втрачається, оскільки $S_A$ містить $c_A$ і $\Pi$, а $Z$ містить маску, координати та відомості про походження.

Часткові оцінки обчислюють до узгодження і зберігають у складі $z$. Метод аналізу спектральних, зорових і структурних ознак формує оцінки $\rho_v$ і $\rho_s$. Метод поведінкового узгодження ознак різного походження формує оцінку ознак подій і поведінки $\rho_b$ й оцінку прояву впливу $Q_A$ з урахуванням причинно пов’язаних подій. Правила обчислення цих величин наведено в розділі~\ref{chap:methods_technology}.

Оцінки $\rho_b$ і $Q_A$ походять з одного представлення ознак подій і поведінки $Z_b$, тому їх не можна вважати незалежними доданками. Перша характеризує виразність нетипових подій, а друга характеризує їх відповідність компонентам і послідовності оброблення. Кон’юнктивне правило їх спільного врахування визначено формулою~\eqref{eq:behavior_joint_score}:
\begin{equation}
\begin{aligned}
F_{bq}(\rho_b,Q_A,C_I)&=\rho_bQ_A,
&& \rho_b,Q_A\in[0,1]\ \text{і походження підтверджено в }C_I;\\
F_{bq}(\rho_b,Q_A,C_I)&=\bot,
&& \text{в іншому разі}.
\end{aligned}
\label{eq:behavior_joint_score}
\end{equation}

де $F_{bq}$ позначає функцію спільного врахування залежних оцінок подій і поведінки. Добуток не трактує їх як окремі свідчення: поведінковий внесок є високим лише за одночасної виразності подій і їх причинно-фазової відповідності досліджуваному компоненту.

Для застосування маски позначимо $r_v=\rho_v$, $r_s=\rho_s$ і $r_{bq}=F_{bq}(\rho_b,Q_A,C_I)$. Маска спільної поведінкової оцінки має значення $m_{bq}=1$, якщо $m_b=1$, значення $r_{bq}$ доступне та його походження підтверджено у $C_I$; в інших випадках $m_{bq}=0$.

Множина фактично доступних часткових оцінок має вигляд $I_{\rho}=\{k\in\{v,s,bq\}:m_k=1\}$. За $I_{\rho}\neq\varnothing$ інтегральну оцінку доказових підстав подано формулою~\eqref{eq:integral_risk}:
\begin{equation}
\rho_A=\sum_{k\in I_{\rho}}\omega_k r_k,
\label{eq:integral_risk}
\end{equation}

де $\rho_v$ і $\rho_s$ є частковими оцінками складової спектральних і зорових ознак та структурної складової; $\rho_b$ позначає оцінку ознак подій і поведінки; $Q_A$ є оцінкою їх зв’язку з проявом прихованого впливу через компонент вебсистеми; $F_{bq}$ позначає функцію спільного врахування залежних оцінок подій і поведінки; $C_I$ містить відомості про походження та конфігурацію інтеграції; $\omega_k$ є ваговим коефіцієнтом доступної оцінки $r_k$; $\rho_A$ позначає інтегральну оцінку доказових підстав. Сумування лише за $I_{\rho}$ усуває невизначену операцію множення нульової ваги на $\bot$.

Доступні часткові оцінки та значення $F_{bq}$ належать відрізку $[0,1]$. Щоб відокремити параметри конфігурації від ваг, одержаних після застосування маски, у $\mathrm T_D$ зберігають базові ваги $\bar\omega_v,\bar\omega_s,\bar\omega_{bq}>0$. Для $k\in\{v,s,bq\}$ вагу з урахуванням маски визначено формулою~\eqref{eq:masked_weights}:
\begin{equation}
\omega_k=\frac{m_k\bar\omega_k}
{m_v\bar\omega_v+m_s\bar\omega_s+m_{bq}\bar\omega_{bq}},\qquad k\in\{v,s,bq\},
\label{eq:masked_weights}
\end{equation}

де $\omega_k$ позначає нормований ваговий коефіцієнт доступної часткової оцінки $r_k$; $m_{bq}$ позначає доступність саме спільної поведінкової оцінки. Відповідно до формули~\eqref{eq:behavior_joint_score}, рівність $m_{bq}=1$ одночасно потребує $m_b=1$, доступних $\rho_b$ і $Q_A$ та підтвердженого в $C_I$ походження; сама наявність журналу або окремого значення $Q_A$ для цього недостатня.

За додатного знаменника формули~\eqref{eq:masked_weights} ваги після застосування маски є невід’ємними та мають одиничну суму, тому $\rho_A\in[0,1]$. Недоступна складова не входить до множини $I_{\rho}$ і не бере участі в сумуванні.

Якщо жодна часткова оцінка недоступна, знаменник дорівнює нулю, $I_{\rho}=\varnothing$, аналітичний стан $S_A$ не формують і реєструють неможливість аналізу. Зокрема, це виконується за $m_A=(0,0,0)$. Одержаний за неповного складу даних результат не вважають рівноцінним результату повної перевірки; неповноту відображають у $c_A$ та поясненні $\Pi$.

Якщо доступна лише одна часткова оцінка, її нормована вага дорівнює одиниці, а $\rho_A$ чисельно збігається з цією оцінкою. Така рівність не перетворює однокомпонентний результат на повний. Вона лише зберігає шкалу доступної складової в межах відповідної маскозалежної конфігурації. У $c_A$ і $\Pi$ при цьому має бути відображено, що висновок не містить підтвердження іншими групами.

За двох або трьох доступних оцінок нормування перерозподіляє ваги лише між фактично наявними складниками. Воно не оцінює якість відсутньої складової й не компенсує її доказовий внесок. Зміна маски змінює знаменник формули~\eqref{eq:masked_weights}, а отже, і функціональний контекст $\rho_A$. Саме тому результати різних масок не порівнюють як значення однієї відкаліброваної шкали без окремої перевірки.

Граничне значення часткової оцінки також не має самостійного предметного висновку. Значення $0$ є виміряним результатом доступної складової, а не позначенням її відсутності. Значення $1$ є верхньою межею визначеної оцінки, але не доводить шкідливої семантики або фактичного впливу. Тлумачення обох меж залежить від походження, маски та правила одержання часткової оцінки.

Правило~\eqref{eq:behavior_joint_score} однозначно визначає спільне врахування $\rho_b$ і $Q_A$, однак у виконаних експериментах його не калібровано через відсутність повного поведінкового корпусу. Так само не визначено числових значень базових ваг і окремого набору для калібрування всіх допустимих масок доступності.

Формула~\eqref{eq:integral_risk} задає виконувану структуру обчислення після фіксації ваг, але не є перевіреною ймовірнісною моделлю. Результати для різних масок можна формально обчислювати за повністю заданої конфігурації $\mathrm T_D$, однак без окремого маскоспецифічного калібрування немає підстав вважати їх рівносильними або порівнювати за єдиною шкалою.

До незалежного добору параметрів і перевірки калібрування за показником Брієра або діаграмою надійності тлумачення результату має спиратися на часткові оцінки, маску доступності та експертний розгляд; результат не використовують як експериментально підтверджене правило реагування.

Індекс класу з найбільшою оцінкою визначають за формулою~\eqref{eq:class_candidate}:
\begin{equation}
j^{*}=\arg\max_{j=1,\ldots,n_A}p_j,
\label{eq:class_candidate}
\end{equation}

де $j^{*}$ позначає індекс найбільшої класової оцінки. Якщо $\rho_A<\tau$, встановлюють $\hat k=\mathrm K_0$ і $\hat y=0$. За умови $\rho_A\geq\tau$ встановлюють $\hat y=1$, але конкретний клас $\mathrm K_{j^{*}}$ призначають лише тоді, коли $p_{j^{*}}\geq\tau_p$ і $p_{j^{*}}-p_{j^{(2)}}\geq\delta_p$, де $p_{j^{(2)}}$ є другою за величиною класовою оцінкою. Якщо хоча б одну з цих двох умов не виконано, встановлюють $\hat k=\mathrm K_U$. Однакові максимальні оцінки впорядковують за правилом, збереженим у $\mathrm T_D$, лише для відтворюваного обчислення $j^{*}$; нульовий розрив між ними не дає підстав приписувати конкретний тип впливу.

Порогове правило розмежовує факт досягнення загального рівня підстав і визначеність конкретного типу. Високе значення окремого $p_j$ не змінює результату $\mathrm K_0$, якщо $\rho_A<\tau$, оскільки сукупні підстави не досягли загального порога. Навпаки, за $\rho_A\geq\tau$ недостатня класова оцінка або малий розрив приводять до $\mathrm K_U$: підстави для захисної дії наявні, але тип не встановлено. Це запобігає примусовому вибору конкретного класу лише через існування найбільшого елемента вектора $p$.

Стан $\mathrm K_U$ не є помилкою обчислення й не означає відсутності результату. Він є визначеним аналітичним станом за достатньої загальної виразності та недостатньої класової визначеності. Помилка або неможливість аналізу виникає раніше, коли немає допустимого $Z$, повної конфігурації $\mathrm T_D$ або жодної доступної часткової оцінки. Розмежування цих випадків не дозволяє підмінити технічну відмову невизначеним класом.

Бінарний висновок призначено для передавання результату до правила реагування. Значення $\hat y=1$ означає, що доступні ознаки досягли встановленого порога для подальшої захисної дії. Значення $\hat y=0$ означає лише недосягнення цього порога. Жодне зі значень не доводить виконання шкідливого коду і не гарантує безпечності об’єкта.

Області визначення та значень функцій упевненості й пояснення задано формулою~\eqref{eq:quality_explanation_domains}:
\begin{equation}
\begin{gathered}
F_q:\mathrm Z\times\mathrm P^{n_A-1}\times\mathrm T_D\rightarrow[0,1],\\
F_E:\mathrm Z\times\mathrm K\times[0,1]^2\times\mathrm T_D\rightarrow\mathrm E,
\end{gathered}
\label{eq:quality_explanation_domains}
\end{equation}

де $\mathrm Z$ позначає множину допустимих інтегрованих представлень; $\mathrm P^{n_A-1}$ є множиною нормованих векторів класових оцінок; $\mathrm K$ позначає множину класів стану; $\mathrm T_D$ є множиною параметрів прийняття рішення; $\mathrm E$ позначає множину структурованих пояснень. Упевненість обчислюють як $c_A=F_q(Z,p;\mathrm T_D)$. Вона залежить від співвідношення класових оцінок, складу доступних даних за $m_A$, їх якості, суперечливості ознак і причинної придатності подій. Низька упевненість не зменшує інтегральної оцінки доказових підстав автоматично. Високе значення $\rho_A$ за низького $c_A$ вказує на значущі, але неповні підстави. Низькі значення обох показників також не підтверджують нормального стану за неповного спостереження.

Поєднання значень $\rho_A$ і $c_A$ тлумачать лише за правилами версованої конфігурації $\mathrm T_D$. За визначеної та перевіреної функції $F_q$ більше значення $c_A$ може відповідати вищій оцінці повноти, якості й несуперечливості підстав у межах цієї конфігурації. Воно не розширює висновок за межі визначених класів і доступних спостережень. Високе $\rho_A$ за нижчого $c_A$ зберігає пороговий результат, але його допустиме тлумачення визначає конкретне правило $F_q$; зовнішня дія при цьому не входить до моделі.

Значення $\rho_A<\tau$ установлює недосягнення порога незалежно від $c_A$, але не гарантує абсолютної безпечності. Значення $c_A$ додатково характеризує підстави цього результату лише відповідно до визначеного $F_q$. Пояснення повинно відображати відмінність маски, якості й суперечливості спостережень, навіть якщо двійкове значення $\hat y$ у кількох випадках однакове. Без перевіреного калібрування $F_q$ значення $c_A$ не використовують для універсального ранжування сили висновків між конфігураціями.

Функція $F_q$ не повинна автоматично дублювати $\rho_A$. Якщо упевненість визначати лише з інтегральної оцінки, неповнота маски, суперечливість і технічна якість перестануть впливати на результат окремо. Контракт $F_q$ вимагає врахування структури $Z$ та вектора $p$ за правилами $\mathrm T_D$. Конкретний числовий вигляд цих правил не задається моделлю, але має бути повністю визначений до формування відтворюваного $S_A$.

Причину рішення відтворює пояснення $\Pi=F_E(Z,\hat k,\rho_A,c_A;\mathrm T_D)$. Для структурного відхилення в ньому зазначають тип і байтовий проміжок, для відхилення у спектральних і зорових ознаках указують просторову та часову ділянку, а для відхилення в ознаках подій і поведінки визначають фазу оброблення і причинне посилання. Якщо рішення ґрунтується на кількох групах ознак, пояснення відображає зв’язок між ними. У ньому також зазначають недоступні складові, причини неповноти, виявлені суперечності та версію параметрів $\mathrm T_D$.

Мінімальний контракт пояснення включає результат, підстави та межі. Результат подають через $\hat y$, $\hat k$, $\rho_A$ і $c_A$. Підстави містять доступні часткові оцінки та посилання на відповідні координати. Межі охоплюють маску, недоступні складові, показники якості, суперечності та версію конфігурації. Якщо хоча б одну з цих частин пропущено, пояснення може описувати значення, але не забезпечує повної перевірки шляху його формування.

Пояснення не повинно приписувати ознаці ширший зміст, ніж установлено відповідною групою. Байтове відхилення описують як структурну підставу, просторово-часове відхилення визначають як спектральну або зорову підставу, а подієве відхилення враховують лише за підтвердженої причинної належності. Висновок про шкідливий код потребує підтвердженої програмної семантики й не виникає автоматично з високого $\rho_A$. Отже, $\Pi$ забезпечує предметну точність аналітичного стану, а не лише перелік найбільших числових складників.

Реалізація функцій $F_D$, $F_q$ і $F_E$ має відповідати трьом перевірюваним вимогам. По-перше, вилучення доступної складової за незмінних інших умов не повинно збільшувати $c_A$. По-друге, якщо всі доступні часткові оцінки однакові, нормування ваг не повинно змінювати $\rho_A$. По-третє, за незмінних $Z$ і $\mathrm T_D$ повторне обчислення має формувати той самий $S_A$; для наскрізного процесу додатково потрібні незмінні $X_A$, правила перетворення, версії засобів і всі параметри.

Для відтворюваного обчислення всі ваги, пороги й правила мають бути повністю задані у версованій конфігурації $\mathrm T_D$. Формула~\eqref{eq:quality_explanation_domains} задає області визначення та значень функцій, але не підміняє їх числової специфікації: конкретний вигляд $F_q$ і правила побудови $F_E$ повинні входити до цієї конфігурації. Виконання зазначених вимог перевіряють для кожної конкретної реалізації.

Вхідні дані функції $F_D$, склад аналітичного стану та межу його використання показано на рис.~\ref{fig:decision_criteria}. Інтегральна оцінка доказових підстав, упевненість і пояснення залишаються окремими складниками $S_A$. Пунктирний зв’язок позначає подальше передавання результату до правила реагування: захисна дія не входить до аналітичного стану, тому його можна повторно перевірити незалежно від застосованої дії.

\begin{figure}[htbp]
\centering
\begin{tikzpicture}[x=1cm,y=1cm]
\node[dissertation block,text width=2.85cm] (joint) at (-5.55,0.35)
  {Представлення\\$Z$};
\node[dissertation block,text width=2.85cm] (decision) at (-2.15,0.35)
  {Функція рішення\\$F_D$};
\node[dissertation block,text width=2.85cm] (state) at (1.25,0.35)
  {Аналітичний стан\\$S_A$};
\node[dissertation block,text width=3.40cm] (classification) at (5.05,2.15)
  {Визначення класу\\$p,\hat y,\hat k$};
\node[dissertation block,text width=3.40cm] (risk) at (5.05,0.35)
  {Ризик і упевненість\\$\rho_A,c_A$};
\node[dissertation block,text width=3.40cm] (explanation) at (5.05,-1.45)
  {Перевірюване пояснення\\$\Pi$};

\draw[dissertation arrow] (joint.east) -- (decision.west);
\draw[dissertation arrow] (decision.east) -- (state.west);
\draw[dissertation arrow] (state.east) -- ++(0.45,0) |- (classification.west);
\draw[dissertation arrow] (state.east) -- (risk.west);
\draw[dissertation arrow] (state.east) -- ++(0.45,0) |- (explanation.west);
\end{tikzpicture}
\caption{Формування аналітичного стану моделі без підміни його технологічною дією}
\label{fig:decision_criteria}
\end{figure}


За повної фіксації функцій і параметрів правило прийняття рішення визначає відтворюваний порядок формального обчислення та зберігає відмінність між виразністю доказових підстав і їх надійністю. Практичне тлумачення порогового рішення та його передавання до правила реагування потребують окремої перевірки фактичної конфігурації й маски доступності. Без такої перевірки сформований результат зберігає статус формальної схеми. Маскоспецифічне калібрування є умовою повної реалізації, а не результатом, установленим у цьому дослідженні.

\section{Процес застосування моделі до виявлення ознак аномального вмісту}
\label{sec:detection_process_model}

Модель визначає склад даних, зв’язки між ознаками та аналітичний результат. Для її застосування необхідний порядок переходів, який не допускає передчасного змішування статичних даних із даними про події та поведінку й не включає захисної дії до $S_A$. Межі між моделлю, двома методами та інформаційною технологією визначають допустимий склад даних і відповідальність кожного етапу; конкретні методи й технологію наведено в розділі~\ref{chap:methods_technology}.

Спочатку метод аналізу спектральних, зорових і структурних ознак формує $Z_v$, $Z_s$, $\rho_v$, $\rho_s$, $L_v$ і $L_s$, з яких утворюється статичний опис $H$. Метод поведінкового узгодження одержує початковий журнал, підтверджує причинну належність подій, формує $X_b$ і $Z_b$, після чого виконує $F_I$. Функція $F_D$ перетворює одержане $Z$ на аналітичний стан $S_A$. Захисну дію визначають після завершення аналізу за правилами інформаційної технології.

Межі між цими етапами є межами відповідальності за дані. Перший метод відповідає за статичні ознаки, їх локалізації, часткові оцінки та походження. Другий метод не змінює вже сформованих статичних складників, а перевіряє події, формує подієво-поведінкове представлення й виконує узгодження. Функція $F_D$ використовує інтегроване представлення, але не виправляє його маску або відомості про походження. Інформаційна технологія приймає завершений аналітичний стан і визначає дію без зміни зафіксованих підстав рішення.

Кожне передавання потребує перевірки вхідного й вихідного контрактів. Для $H$ перевіряють склад статичних ознак, відповідність $m_c$, наявність координат і версій схем. Для приєднання $X_b$ перевіряють ідентифікатор запуску, часові межі, фазову належність і повноту журналу. Для $Z$ перевіряють узгодженість маски $m_A$, відомостей $C_I$ і фактично доступних складників. Для $S_A$ перевіряють повноту $\mathrm T_D$ та відтворюваність усіх його складників. Вихід попереднього етапу не передають далі лише на підставі технічного завершення процедури.

Неприпустима підміна відповідальності між етапами. Відсутній $Z_v$ не відновлюють у методі поведінкового узгодження за подіями, а неповний журнал не компенсують високою статичною оцінкою. Функція рішення не створює координат, якщо їх не сформовано в ознаковому ядрі. Правило реагування не змінює клас або упевненість, щоб узгодити їх із бажаною дією. Така незмінність переходів зберігає можливість установити, на якому етапі виникло кожне значення.

Межею завершення аналізу є перехід $S_A=F_D(Z,\mathrm T_D)$, визначений формулою~\eqref{eq:analytic_state}. До його виконання перевіряють походження представлень, сумісність схем ознак, правильність маски доступності, а також повноту й версію конфігурації $\mathrm T_D$.

За відсутності повністю заданої й версованої $\mathrm T_D$ аналітичний стан $S_A$ не формують і реєструють неможливість відтворюваного обчислення. За повної $\mathrm T_D$ та допустимо неповної маски функція $F_D$ може сформувати $S_A$, але неповноту спостереження відображають у $c_A$ та поясненні $\Pi$.

За непідтвердженого походження, несумісності схем ознак або некоректної маски $S_A$ також не формують. Дані іншої перевірки для заповнення відсутніх складників не використовують.

Передумови формування $S_A$ поділяють на структурні та параметричні. Структурні передумови стосуються допустимого $Z$: складники мають належати одному входу, відповідати масці та мати перевірюване походження. Параметричні передумови стосуються $\mathrm T_D$: пороги, базові ваги, правила класового вибору, упевненості й пояснення повинні бути задані та версовані. Виконання лише однієї групи передумов недостатнє. Допустимий $Z$ без повної $\mathrm T_D$ не дає відтворюваного рішення, а повна $\mathrm T_D$ не може виправити недопустимий інтегрований вхід.

Результат проходження цієї межі має три основні статуси. Повний аналітичний стан формують за повного допустимого складу й заданої конфігурації. Явно неповний стан формують за маски, яку підтримує конфігурація, зі збереженням причин неповноти в $c_A$ та $\Pi$. Повідомлення про неможливість аналізу формують, коли немає мінімального складу даних, порушено походження або конфігурація не визначає потрібного обчислення. Останній статус не містить $\hat y=0$, оскільки неможливість сформувати рішення не є недосягненням порога.

Незмінність предметного тлумачення між етапами забезпечують шість правил:

\begin{enumerate}
\item \textit{Єдність входу.} Усі представлення однієї перевірки походять від того самого $X_{mm}$.
\item \textit{Незмінність початкового об’єкта.} Жодне похідне представлення не заміщує $X_{mm}$.
\item \textit{Розділення даних.} Статичний опис не містить представлення ознак подій і поведінки до етапу узгодженого об’єднання.
\item \textit{Явне позначення доступності.} Відсутнє представлення не тлумачать як нульову ознаку.
\item \textit{Збереження походження.} Узагальнена ознака має зв’язок із байтовою, просторово-часовою або подієвою координатою.
\item \textit{Розмежування рішення та дії.} Стан $S_A$ є аналітичним результатом, а захисну дію визначає інформаційна технологія.
\end{enumerate}

Єдність входу й незмінність початкового об’єкта забезпечують тотожність предмета перевірки. Розділення даних та явне позначення доступності визначають, які спостереження фактично брали участь у кожному обчисленні. Збереження походження забезпечує можливість повернутися від результату до відповідної ділянки або події. Розмежування рішення та дії не допускає зміни наукового тлумачення через зовнішню політику. Разом ці правила утворюють цілісний контракт передавання, у якому зміна будь-якого складника повинна бути явно зареєстрована.

Контракт не вимагає, щоб усі перевірки завершувалися повним результатом. Його вимога полягає в однозначності статусу кожного переходу. Якщо складник доступний, мають бути відомі його значення, походження й якість. Якщо він недоступний, мають бути відомі причина та етап виникнення обмеження. Якщо перехід неможливий, наступний етап не повинен створювати результат, який приховує цю відмову. Така організація забезпечує однакове предметне тлумачення штатних, неповних і відмовних проходжень.

Кожний метод завершується повним результатом, явно неповним результатом або повідомленням про неможливість сформувати мінімально необхідний вихід. Неповний результат передають далі разом із маскою доступності, показником якості та причиною обмеження. Якщо походження мінімально необхідного виходу не підтверджено, аналітичний висновок не формують. Такий стан не можна тлумачити як відповідність нормі.

Повний і неповний результати відрізняються складом підстав, але обидва повинні мати відтворюване походження. Неповнота є допустимим модельним станом, якщо її передбачено контрактом відповідної маски й явно відображено в супровідних відомостях. Після технічної помилки частковий результат можна передати далі лише тоді, коли фактично наявні складники самостійно відповідають такому контракту; в іншому разі реєструють неможливість виходу. Це правило визначає допустимість даних, але не задає алгоритму оброблення технічної помилки.

У супровідних відомостях причину обмеження пов’язують з етапом, на якому її встановлено. Помилка розкодування стосується формування $X_v$, несумісність схеми виникає під час формування або передавання ознак, неповне фазове покриття стосується відбору $X_b$, а непідтримувана маска вказує на інтеграцію або прийняття рішення. Така локалізація не визначає процедури відновлення, але дає змогу встановити, які підстави залишилися придатними та який контракт не виконано.

Модель не встановлює числових значень порогів, ваг та інших параметрів $\mathrm T_D$. Їх визначають на даних для налаштування і перевіряють на відокремлених даних. За відсутності такої перевірки формули описують порядок обчислення, але не підтверджують переваги конкретної конфігурації чи її готовності до промислового застосування. Вихід також не вважають коректним, якщо неможливо відтворити його походження або версію схеми ознак.

Версування $\mathrm T_D$ охоплює всі елементи, включені до її визначення: пороги, базові ваги, порядок розв’язання однакових класових оцінок і версії правил обчислення упевненості та пояснення. Збережена версія також має однозначно визначати область застосування конфігурації; якщо певна маска до неї не належить, $S_A$ не формують. Зміна будь-якого з визначених елементів утворює іншу конфігурацію прийняття рішення. Результати різних версій можна зіставляти лише із зазначенням відмінностей, оскільки однаковий $Z$ може породити різні $S_A$ за різних правил.

Повнота конфігурації є передумовою формального обчислення, а її перевіреність є передумовою практичного порогового тлумачення. Повністю задана $\mathrm T_D$ дає змогу відтворити результат, але сама по собі не підтверджує правильності вибраних порогів або ваг. Для цього потрібне окреме оцінювання на відповідних даних. Модель зберігає цю відмінність і не підміняє визначеність процедури доказом її кількісної переваги.

Застосування моделі до нового мультимедійного формату потребує незмінного $X_{mm}$, контрольовано одержаних $X_v$, $X_f$, $X_s$, $X_m$ і, для відео, $X_t$, а також правильних значень $m_A$ та правил визначення координат. Виконання цих умов створює можливість розширення моделі, але не підтверджує підтримки формату без реалізованого засобу аналізу й експериментальної перевірки. Якщо формат допускає кілька коректних тлумачень, у $C$ зазначають застосоване правило, а у висновку визначають межі його чинності.

Розширення на формат починається не з повторного використання наявного числового вектора, а з визначення контрактів представлень. Потрібно встановити, які байтові елементи утворюють $X_s$ і $X_m$, як контрольовано одержують $X_v$, чи застосовні $X_f$ і $X_t$, а також як координати похідних даних пов’язуються з початковим носієм. Лише після цього можна перевіряти сумісність зі схемами $Z_s$ і $Z_v$. Збіг розмірності без такого предметного відображення не підтверджує підтримки формату.

Кілька коректних тлумачень того самого контейнера утворюють кілька похідних спостережень із власними параметрами, а не один довільно вибраний «правильний» результат. Якщо модельна конфігурація підтримує лише одне тлумачення, цей вибір фіксують у $C$ як межу застосування. Розбіжність результатів різних операторів може бути предметом окремого зіставлення, але її не подають як аномалію без визначеного правила. Такий підхід відокремлює неоднозначність формату від прихованого навантаження.

Непідтримуваний формат, невдала реалізація аналізатора й недоступне окреме представлення також мають різні статуси. У першому випадку немає визначеного контракту для формату. У другому контракт існує, але конкретне виконання не сформувало придатного результату. У третьому решта представлень може залишатися доступною за допустимої маски. Точне розмежування визначає, чи можливий частковий аналіз, повторення операції або лише завершення з повідомленням про відсутність підтримки.

На рис.~\ref{fig:detection_process} показано зв’язок чотирьох складників: моделі, двох методів та інформаційної технології. Модель визначає склад і зміст аналітичного стану, методи конкретизують обчислення, а інформаційна технологія встановлює порядок виконання та реагування.

\begin{figure}[htbp]
\centering
\begin{tikzpicture}[x=1cm,y=1cm]
\node[dissertation block,text width=9.8cm,minimum height=12mm] (model) at (0,3.6)
  {Модель мультимодального представлення\\прихованої загрози};

\node[dissertation block,text width=5.2cm,minimum height=13mm] (method1) at (-3.6,1.2)
  {Метод аналізу\\спектральних, зорових\\і структурних ознак};
\node[dissertation block,text width=5.2cm,minimum height=13mm] (method2) at (3.6,1.2)
  {Метод поведінкового\\узгодження ознак\\різного походження};

\node[dissertation block,text width=9.8cm,minimum height=12mm] (technology) at (0,-1.4)
  {П’ятиетапна інформаційна технологія\\виявлення ознак аномального вмісту};

\draw[dissertation arrow] (model.south) -- ++(0,-0.35) -| (method1.north);
\draw[dissertation arrow] (model.south) -- ++(0,-0.35) -| (method2.north);
\draw[dissertation arrow] (method1.east) -- (method2.west);
\draw[dissertation arrow] (method1.south) -- ++(0,-0.35) -| (technology.north);
\draw[dissertation arrow] (method2.south) -- ++(0,-0.35) -| (technology.north);
\end{tikzpicture}
\caption{Взаємозв’язок моделі, двох методів та інформаційної технології}
\label{fig:detection_process}
\end{figure}


Процес застосування визначає послідовність переходів, межі довіри та стани, за яких результати можна узгоджувати й передавати до правила реагування. Два методи та п’ятиетапна інформаційна технологія конкретизують ці положення без зміни змісту модельних виходів у розділі~\ref{chap:methods_technology}.

\section*{Висновки до розділу 2}
\phantomsection
\addcontentsline{toc}{section}{Висновки до розділу 2}

У розділі розроблено модель мультимодального представлення прихованої загрози як формальне розв’язання наукової суперечності, що полягає в багаторівневому прояві прихованої загрози за переважно ізольованого формування локальних оцінок. Модель охоплює умови прихованої доставки, багаторівневе представлення об’єкта, формування ознак, критерії рішення та порядок застосування. Предметне тлумачення моделі ґрунтується на розмежуванні аномалії, аномального коду, прихованої загрози, шкідливого коду, прихованого навантаження та прихованого впливу.

Наукова відмінність моделі полягає у типізованому поєднанні незмінного байтового подання, похідних представлень, локалізацій, маски доступності та відомостей про походження в межах однієї перевірки, а також у відокремленні аналітичного стану від технологічної дії. Початковий об’єкт $X_{mm}$, похідні представлення $X_v$, $X_f$, $X_s$, $X_m$, $X_t$, причинно пов’язані події $X_b$ та ознаки $Z_v$, $Z_s$, $Z_b$ поєднано спільним ідентифікатором перевірки.

Маска $m_A$ відрізняє відсутнє спостереження від виміряного нульового значення, а координати й відомості $C$ зберігають якість та походження даних. Доказовими опорами відмінності є встановлена в розділі 1 прогалина щодо роз’єднаного подання ознак, формальний контракт моделі та перевірка доступних статичних складників.

Аналітичний стан $S_A$ містить оцінки належності до класів, двійковий висновок, визначений клас, інтегральну оцінку доказових підстав, упевненість і пояснення. Залежні оцінки подій і поведінки $\rho_b$ і $Q_A$ враховують спільно через $F_{bq}$, що формально виключає їх подання як незалежних доданків $\rho_A$. Інтегральну оцінку доказових підстав відокремлено від упевненості, тому неповнота й суперечливість підстав не приховуються пороговим рішенням. Захисну дію також відокремлено від $S_A$ і віднесено до інформаційної технології.

Межі одержаного результату визначаються підтримуваними форматами, якістю контрольованого розкодування, повнотою причинно пов’язаних подій і перевіреністю параметрів $\mathrm T_D$.

Повне інтегроване представлення $Z$ із причинно пов’язаними подіями експериментально не формувалося; кількісну перевагу повного мультимодального контуру не підтверджено. Для $F_{bq}$ задано консервативне кон’юнктивне правило, однак у наявних матеріалах не зафіксовано ваг інтегральної оцінки та не проведено калібрування для всіх масок доступності. Тому результати різних масок не вважають рівносильними й не порівнюють за єдиною шкалою без окремої перевірки.

Модель визначає формальні умови аналізу, але не підтверджує практичного порогового тлумачення неперевіреної конфігурації, не гарантує виявлення довільного прихованого навантаження й не підтверджує промислової готовності. Два методи та інформаційна технологія конкретизують модель у розділі~\ref{chap:methods_technology} у межах зазначених доказових обмежень.
